\section{Local trace synthesis and its geometric cost}
\label{sec:constructive}

The geometric multiplier gives a target load, not independent acoustic Cauchy
data. Nevertheless, a conditional construction clarifies what very rich source
access would permit. This section is a deduction for the inviscid transmission
reduction; it is not an experimental fabrication result.

\subsection{Two-fluid traces and a pressure node}
Pressure and normal velocity are continuous at the ideal interface. Write their
traces as $P$ and $v$, and $\Delta\rho=\rho_- -\rho_+$. Eliminating tangential
velocities from \cref{eq:radiation-flux} gives
\begin{equation}
 \Pi=\underbrace{\frac14\left(\frac1{\rho_-c_-^2}
                       -\frac1{\rho_+c_+^2}\right)}_{c_P}|P|^2
       +\frac{\Delta\rho}{4}|v|^2
       +\frac{\Delta\rho}{4\omega^2\rho_-\rho_+}|\Gs P|^2.
 \label{eq:trace-traction-constructive}
\end{equation}
The tangential radiation-traction jump vanishes in this ideal model. The
coefficients and orientation follow from the momentum-flux convention, as in
the coupled-wave approach of \cite{chesneau2022}.

Set $H=\sigma\kappa+\Delta\rho\Phi$. If $\Delta\rho\ne0$ and an admissible
pressure offset makes $(H-\lambda_V)/\Delta\rho>0$, then
\begin{equation}
 P|_\Gamma=0,\qquad v=2\sqrt{(H-\lambda_V)/\Delta\rho}
 \label{eq:pressure-node-traces}
\end{equation}
realizes the required normal load at the trace level. If $\Delta\rho=0$ but
$c_P\ne0$, the alternative $v=0$, $P=\sqrt{(H-\lambda_V)/c_P}$ does so.
Equal density and sound speed give zero ideal radiation-traction contrast,
irrespective of optical-index contrast. These trace prescriptions have yet to
be produced by admissible sources.

\subsection{Conditional local fields and volume forces}
For analytic target patches and analytic positive radicand, let
$\varphi_\ell$ solve $(\Delta+k_\ell^2)\varphi_\ell=0$ locally with
$\varphi_\ell=0$, $\partial_n\varphi_\ell=v$ on the patch and
$k_\ell=\omega/c_\ell$. The analytic Cauchy theorem \cite{kowalevsky1875} applies to this
noncharacteristic surface. It is a local existence statement; elliptic Cauchy
continuation to remote walls need not be stable. Let smooth cutoffs $\chi_\ell$
equal one near the patch and vanish inside the available analytic neighborhood.
With $\widetilde\varphi_\ell=\chi_\ell\varphi_\ell$, define
\begin{align}
 \bm V_\ell&=\nabla\widetilde\varphi_\ell,&
 P_\ell&=\frac{\rho_\ell c_\ell^2}{\mathrm i\omega}
                              \Delta\widetilde\varphi_\ell,\\
 \bm F_\ell&=\frac{\rho_\ell c_\ell^2}{\mathrm i\omega}
                   \nabla(\Delta+k_\ell^2)\widetilde\varphi_\ell,&q_\ell&=0.
 \label{eq:cutoff-volume-force}
\end{align}
Substitution verifies both forced first-order acoustic balances. The force is
supported where the cutoff changes, leaving the prescribed interface traces
unaltered. This requires independently realizable volume forces on both sides,
with compatible mean reactions and energy budgets. At contact lines and patch
edges the construction requires additional matching. For two separated faces,
disjoint source neighborhoods may supply local constructions, but a single
boundary-only Helmholtz problem imposes additional joint compatibility.
Neither local construction proves controllability from a cylindrical array.

\subsection{A necessary bound for strict geometric accuracy}
Let $\epsilon_a$ bound peak normal carrier displacement. Since
$v=-\mathrm i\omega\Xi$, a pressure-node realization obeys
$|\Pi|\leq |\Delta\rho|\omega^2\epsilon_a^2/4$.
Its load has one sign. Consequently a necessary condition, independent of the
permitted pressure offset, is
\begin{equation}
 \operatorname{osc}_{\Gamma}H
 \leq\frac{|\Delta\rho|\omega^2\epsilon_a^2}{4},\qquad
 f\geq\frac{1}{\pi\epsilon_a}
            \sqrt{\frac{\operatorname{osc}_{\Gamma}H}{|\Delta\rho|}}.
 \label{eq:pressure-node-displacement-bound}
\end{equation}
Here $\operatorname{osc}H=\sup H-\inf H$. A strictly positive load margin
increases this cost. For graph-height error the vertical oscillation at fixed
horizontal position is $\Xi/n_z$, so the pointwise height bound is stronger
where $|n_z|<1$. This inequality rejects a particular realization when violated;
pressure-dominated or tangential-field realizations may avoid it. Raising
frequency also changes attenuation, thermal loads, acoustic resolution and
actuator response, so it is not by itself a demonstrated solution.
